\section{Research Problem}

This section presents the first step of the Design Science Research Methodology (DSRM), the identification of the problem, and its motivation. 

As shown in the results in Tables 5, 6, 7, 8 and 9 to the research questions, a lot of sectors can take advantage of the AI and LLMs transformative assets, incorporate it into their workflow. Although AI has been in the Software Development landscape for a while now, due to the highly dynamic and fast-evolving nature of the field, new models and new tools are being created nearly every month, and AI usage in general has never been so prevalent not only in the Software Development field but in other industry sectors as well. 

The conclusions drawn from the information collected in the systematic literature review, indicate that while many benefits can be attained from the use of AI, many challenges and disadvantages are in the way of its successful implementation and usage. Therefore the real main challenge organizations face and one of the most important consequences that is mentioned within the literature is the lack of established, solid and clear methodologies, that shape the proper use of AI and LLMs within the Software Development realm and across the industry and professional practice. \cite{Banh:2025cf, Farhanna:2025eu, Hemmat:2025rd, Gao:2025cc}

Although several frameworks and approaches that incorporate AI into the development process have been proposed, as summarized in Table 7, these solutions differ significantly in scope, assumptions, and level of abstraction. They often lack a shared foundation regarding lifecycle integration, management practices, and evaluation criteria. Consequently, organizations face difficulties not only in adopting AI-supported development practices, but also in objectively assessing their impact on software quality, security, and maintainability due to the lack of standardized benchmarks and performance metrics; ultimately proving that the field is actively searching for structured and repeatable integration patterns and methodologies. \cite{Ahmed:2025ai}

This problem is further amplified by the fact that AI increasingly influences all phases of the SDLC. Due to this problem, companies are seeking for a new methodological evolution to fully embrace, utilize, incorporate and implement AI and LLMs into their workflow. Traditional methodologies such as Waterfall, Scrum, Kanban, and the Agile Manifesto introduced substantial improvements in adaptability and collaboration, but they were not designed for development environments in which AI or LLMs actively contribute to software creation. \cite{Zhang:2025ea}

As a result and as AI-assisted development becomes more prevalent, modern development teams experience new dynamics, including the rise of Vibe Coding, accelerated implementation, emerging bottlenecks created by AI in earlier phases such as requirements clarification and later phases such as testing, and a stronger emphasis on delivering product value rather than merely increasing feature output. Together, these changes indicate that existing methodologies may no longer fully align with emerging AI-augmented development dynamics. Existing traditional frameworks often fall short in AI-driven contexts because:

\begin{itemize}
    \item high-speed AI-assisted coding accelerates implementation but creates pressure upstream (requirements clarification) and downstream (testing, verification, and maintenance);
    \item they lack explicit structures for AI governance, prompt engineering, model reliability assessment, and risk mitigation;
    \item they assume all contributors are humans, whereas AI models introduce a new category of collaborator.
\end{itemize}

\textbf{Therefore, the core research problem addressed in this work is the lack of patterned guidance, strict methodologies and standardized evaluation metrics that can be defined through a structured, process-oriented framework that supports the systematic integration and management of AI and LLM-based tools across the Software Development Lifecycle, while preserving human oversight, software quality and security, and organizational accountability. This gap limits organizations’ ability to fully and sustainably realize the potential benefits of AI-assisted software development.}

